%% MIT Thesis class sample appendix with a long table
%% version 1.04, 2025/11/02
\chapter{Derivation of the Saha equation of nuclear statistical equilibrium}\label{ap:Detailed_balance}

This appendix derives the Saha equation, Eq.~\ref{eq:saha}, used in Sec.~\ref{sec:saha}.
It is obtained by assuming a Maxwell-Boltzmann distribution and then using the equilibrium condition $\mu_{Z,A}=Z\mu_{p}+N\mu_{n}$ (Secs.~\ref{sec:apB_ideal}--\ref{sec:apB_abundances}).
%------------------------------------------------------------------------------------------
\section{Ideal non-degenerate gas of particles with internal structure}\label{sec:apB_ideal}
%------------------------------------------------------------------------------------------

Consider species $i$ of mass $m_i$ in a volume $V$ at temperature $T$.
A particle of this species is labeled by its momentum $\mathbf p$ and by its internal state $l$.
The state $l$ has excitation energy $E_l$ and spin $J_l$, and $l=0$ is the ground state.
In the non-relativistic limit the single-particle energy, including the rest mass, is
\begin{equation}\label{eq:apB_energy}
  \varepsilon_{\mathbf p,l}=m_ic^2+E_l+\frac{p^2}{2m_i},
\end{equation}
and each level has the magnetic degeneracy $2J_l+1$.
In the grand-canonical ensemble with chemical potential $\mu_i$, which here includes the rest mass, the mean occupation of a single-particle state is
\begin{equation}\label{eq:apB_occupation}
  \langle n_{\mathbf p,l}\rangle=\frac{1}{\exp\left[(\varepsilon_{\mathbf p,l}-\mu_i)/k_BT\right]\pm1},
\end{equation}
with the upper sign for fermions and the lower for bosons.
If $\exp[(\varepsilon_{\mathbf p,l}-\mu_i)/k_BT]\gg1$, both statistics reduce to the Maxwell--Boltzmann, 
\begin{equation}
    \langle n_{\mathbf p,l}\rangle\simeq\exp[(\mu_i-\varepsilon_{\mathbf p,l})/k_BT].
\end{equation}

\noindent The number density is the sum over states times the occupation of the state.

\begin{align}
    n_i&=\int \frac{d^3\mathbf P}{(2\pi\hbar)^3}\sum_{l}(2J_l+1)\langle n_{\mathbf p,l}\rangle,\\
    n_i&=\int_0^\infty \frac{4\pi p^2dp}{(2\pi\hbar)^3}\sum_{l}(2J_l+1)\exp\left[\frac{1}{k_BT}(\mu_i-m_ic^2-E_l-\frac{p^2}{2m_i})\right],
\end{align}

\begin{equation}\label{eq:apB_n_sum}
  n_i=\exp\left[\frac{1}{k_BT}(\mu_i-m_ic^2)\right]\left[\sum_l(2J_l+1)\,\exp\left(-\frac{E_l}{k_BT}\right)\right]\int_0^\infty dp\frac{ p^2}{2\pi^2\hbar^3}\,\exp\left[-\frac{p^2}{2m_ik_BT}\right].
\end{equation}
The momentum integral is a Gaussian,
\begin{equation}
    \int_0^\infty dp\frac{ p^2}{2\pi^2\hbar^3}\,\exp\left[-\frac{p^2}{2m_ik_BT}\right]=\left[\frac{m_ik_BT}{2\pi \hbar^2}\right]^{\frac{3}{2}}.
\end{equation}
From that we define the thermal wavelength of Sec.~\ref{sec:saha}.
\begin{equation}\label{eq:apB_gauss}
  \left(\frac{m_ik_BT}{2\pi\hbar^2}\right)^{3/2}=\left(\frac{1}{\Lambda_i}\right)^{3},
\end{equation}

\noindent The internal sum is $g_iG_i(T)$, with $g_i=2J_0+1$ and $G_i$ the partition function of Eq.~\ref{eq:partition}.
Therefore
\begin{equation}\label{eq:apB_MB}
  n_i=\frac{g_iG_i(T)}{\Lambda_i^3}\,\exp\left(\frac{\mu_i-m_ic^2}{k_BT}\right),
\end{equation}
which is Eq.~\ref{eq:MB_density}.
Its inverse is the chemical potential of Eq.~\ref{eq:mu_MB},
\begin{equation}\label{eq:apB_mu}
  \mu_i=m_ic^2+k_BT\ln\left(\frac{n_i\Lambda_i^3}{g_iG_i(T)}\right) .
\end{equation}

%-------------------------------------------------------------
\section{The Saha equation for number densities}\label{sec:apB_saha_n}
%------------------------------------------------------------------------------------------

Inserting Eq.~\ref{eq:apB_mu} into both sides of $\mu_{Z,A}=Z\mu_p+N\mu_n$, with $g_p=g_n=2$ and $G_p=G_n=1$ for the nucleons (spin-$\tfrac12$ particles without bound excited states), gives
\begin{multline}\label{eq:apB_step1}
  m(Z,A)c^2+k_BT\ln\frac{n(Z,A)\Lambda_{(Z,A)}^3}{g(Z,A)G(Z,A)}\\
  =Zm_pc^2+Nm_nc^2+k_BT\left[Z\ln\frac{n_p\Lambda_p^3}{2}+N\ln\frac{n_n\Lambda_n^3}{2}\right].
\end{multline}
The rest masses combine into the binding energy of Eq.~\ref{eq:binding}, $B(Z,A)=[Zm_p+Nm_n-m(Z,A)]c^2$.
Dividing by $k_BT$ and exponentiating,
\begin{equation}\label{eq:apB_saha_n}
  n(Z,A)=\frac{g(Z,A)\,G(Z,A)}{\Lambda_{(Z,A)}^3}\left(\frac{n_p\Lambda_p^3}{2}\right)^{Z}\left(\frac{n_n\Lambda_n^3}{2}\right)^{N}\exp\left(\frac{B(Z,A)}{k_BT}\right),
\end{equation}
which is Eq.~\ref{eq:saha_n}.
The masses in $B(Z,A)$ are nuclear masses, because in NSE the atoms are completely ionized and the electrons are counted separately.

%------------------------------------------------------------------------------------------
\section{From number densities to abundances}\label{sec:apB_abundances}
%------------------------------------------------------------------------------------------

All thermal wavelengths are now expressed in terms of $\Lambda_u$, the thermal wavelength at the atomic mass unit.
Exactly,
\begin{equation}\label{eq:apB_lambdas}
  \Lambda_p^3=c_p\,\Lambda_u^3,\qquad \Lambda_n^3=c_n\,\Lambda_u^3,\qquad \Lambda_{(Z,A)}^{-3}=\epsilon_A\,A^{3/2}\,\Lambda_u^{-3},
\end{equation}
with
\begin{equation}\label{eq:apB_mass_factors}
  c_{p}=\left(\frac{m_u}{m_{p}}\right)^{3/2}=0.9892,\qquad c_{n}=\left(\frac{m_u}{m_{n}}\right)^{3/2}=0.9871,\qquad \epsilon_A=\left(\frac{m(Z,A)}{Am_u}\right)^{3/2}.
\end{equation}
Substituting Eq.~\ref{eq:apB_lambdas} into Eq.~\ref{eq:apB_saha_n} and writing $n_i=n_bY_i$,
\begin{equation}\label{eq:apB_exact}
  n_bY(Z,A)=\epsilon_A\,g\,G\,\frac{A^{3/2}}{\Lambda_u^{3}}\,\frac{(n_b\Lambda_u^{3})^{A}}{2^A}\,(c_pY_p)^{Z}\,(c_nY_n)^{N}\exp\left(\frac{B}{k_BT}\right).
\end{equation}
Dividing by $n_b$ and collecting $(n_b\Lambda_u^3)^{A-1}=\theta^{A-1}$, with $\theta$ defined in Eq.~\ref{eq:theta}, yields
\begin{equation}\label{eq:apB_saha_exact}
  Y(Z,A)=\epsilon_A\;g(Z,A)\,G(Z,A)\,\frac{A^{3/2}}{2^A}\,\theta^{A-1}\exp\left(\frac{B(Z,A)}{k_BT}\right)(c_pY_p)^{Z}\,(c_nY_n)^{N}.
\end{equation}
With $\epsilon_A=c_p=c_n=1$ this is Eq.~\ref{eq:saha}.


